\section*{Contexte général}
\addcontentsline{toc}{section}{Contexte général}

L'accès à une information fiable est un besoin important dans le domaine de la santé. Pour les
pharmacies, cette information concerne souvent l'adresse, le numéro de téléphone, les horaires ou encore
les gardes disponibles. Dans une situation urgente, l'utilisateur cherche surtout une réponse rapide :
trouver une pharmacie proche ou savoir quelle pharmacie assure la garde à une date donnée.

En Tunisie, les informations sur les pharmacies de garde existent, mais elles ne sont pas toujours
faciles à consulter. Elles peuvent être publiées sous forme de listes, partagées localement ou mises à jour
manuellement. Cette organisation rend la recherche moins pratique, surtout depuis un téléphone ou dans
une zone que l'utilisateur ne connaît pas bien.

C'est dans ce contexte que s'inscrit le projet \textbf{PharmacieConnect}. La plateforme regroupe une API
backend développée avec FastAPI, un tableau de bord administrateur réalisé avec React/Vite et une
application mobile Expo/React Native. L'objectif est de faciliter l'accès aux données côté citoyen tout en
donnant aux administrateurs un outil de gestion plus structuré.

\section*{Problématique}
\addcontentsline{toc}{section}{Problématique}

La question principale posée par ce projet est la suivante : comment mettre en place une plateforme qui
permet au citoyen de rechercher rapidement une pharmacie ou une pharmacie de garde, tout en permettant
à l'administrateur de gérer les données de manière sécurisée et organisée ?

Cette question implique plusieurs points. L'application mobile doit rester simple à utiliser et proposer
des résultats clairs. Le tableau de bord doit permettre d'importer, de corriger et de suivre les données.
Le backend doit, de son côté, contrôler les accès, valider les informations reçues et garder une trace des
actions importantes.

\section*{Objectifs du projet}
\addcontentsline{toc}{section}{Objectifs du projet}

Le projet a pour objectif de réaliser une plateforme complète de gestion des pharmacies de garde. Les
objectifs retenus sont les suivants :

\begin{itemize}
    \item développer une API backend qui centralise les données et les règles métier ;
    \item mettre en place une authentification sécurisée avec des jetons JWT ;
    \item créer une interface web pour gérer les pharmacies, les gardes, les médicaments et les utilisateurs autorisés ;
    \item développer une application mobile pour la recherche de pharmacies, la consultation des gardes et l'accès au catalogue de médicaments ;
    \item intégrer une recherche par proximité à partir des coordonnées GPS ;
    \item permettre l'import de fichiers CSV avec validation des données et retour d'erreurs ;
    \item ajouter des journaux d'audit, des indicateurs et un contrôle d'accès basé sur les rôles ;
    \item tester les fonctionnalités principales à travers des tests backend et des scénarios fonctionnels.
\end{itemize}

\section*{Méthodologie adoptée}
\addcontentsline{toc}{section}{Méthodologie adoptée}

La méthodologie adoptée pour la réalisation de ce projet est [À compléter]. Le travail a été mené par
étapes : analyse du besoin, étude de l'existant, définition des fonctionnalités, conception de
l'architecture, développement du backend, développement de l'administration web, développement de
l'application mobile, intégration, tests et rédaction du rapport.

Cette progression a permis de commencer par les règles métier du backend avant de les exploiter dans les
interfaces. Elle a aussi facilité la validation des imports, de l'authentification, de la gestion des rôles
et des principaux parcours utilisateur.

\section*{Organisation du rapport}
\addcontentsline{toc}{section}{Organisation du rapport}

Ce rapport est composé de quatre chapitres. Le premier chapitre présente le cadre général du projet, le
domaine pharmaceutique, l'étude de l'existant, ses limites et la solution proposée. Le deuxième chapitre
décrit les acteurs, les besoins fonctionnels, les besoins non fonctionnels et les règles de gestion. Le
troisième chapitre présente la conception de la solution, son architecture, son modèle de données, ses
diagrammes UML et ses mécanismes de sécurité. Le quatrième chapitre détaille la réalisation technique,
l'intégration API, la configuration, les tests et les difficultés rencontrées. Le rapport se termine par une
conclusion générale et des perspectives d'amélioration.
